\begin{examplebox}
	\textbf{\textcolor{CExample}{例 1（基础）}}

	\textbf{\textcolor{CExample}{题目：}}
	在二面角 $\alpha\text{-}l\text{-}\beta$ 中，$OA\perp l$，$OB\perp l$，且 $OA=3$，$OB=4$，$AB=5$。求二面角 $\theta$ 的大小。

	\textbf{\textcolor{CExample}{解题步骤：}}
	\begin{description}[style=nextline, leftmargin=2.8em, labelsep=0.5em, itemsep=0.45em]
		\item[\textbf{第一步（确定平面角）：}] 由定义，$\angle AOB$ 是二面角的平面角，故 $\theta = \angle AOB$。
			\par\textbf{理由：}
			两条射线均垂直于棱，且从同一点出发，满足平面角定义。

		\item[\textbf{第二步（余弦定理代入）：}] 在 $\triangle AOB$ 中，由余弦定理得
			\[
				\cos\theta = \frac{3^2+4^2-5^2}{2\times3\times4}.
			\]
			\par\textbf{理由：}
			直接代入 $OA=3, OB=4, AB=5$。

		\item[\textbf{第三步（计算求角）：}] 计算分子 $9+16-25=0$，分母 $24$，故 $\cos\theta=0$，因此 $\theta=90^\circ$。
			\par\textbf{理由：}
			在 $[0^\circ,180^\circ]$ 内，余弦为 $0$ 对应 $90^\circ$（直二面角）。
	\end{description}

	\textbf{\textcolor{CExample}{关键结论：}}
	\textbf{$\theta=90^\circ$，即二面角为直二面角。}

	\medskip
	\textbf{\textcolor{CExample}{例 2（稍变形）}}

	\textbf{\textcolor{CExample}{题目：}}
	在二面角 $\alpha\text{-}l\text{-}\beta$ 中，$OA\perp l$，$OB\perp l$，已知 $OA=2$，$OB=3$，$AB=\sqrt{7}$，求 $\cos\theta$ 及 $\theta$。

	\textbf{\textcolor{CExample}{解题步骤：}}
	\begin{description}[style=nextline, leftmargin=2.8em, labelsep=0.5em, itemsep=0.45em]
		\item[\textbf{第一步（平面角认定）：}] 同上，$\theta=\angle AOB$。
			\par\textbf{理由：}
			定义法。

		\item[\textbf{第二步（代入公式）：}]
			\[
				\begin{aligned}
					\cos\theta &= \frac{2^2+3^2-(\sqrt{7})^2}{2\times2\times3} \\
					&= \frac{4+9-7}{12} \\
					&= \frac{6}{12} \\
					&= \frac12.
			\end{aligned}
			\]
			\par\textbf{理由：}
			注意 $(\sqrt{7})^2=7$，逐步化简。

		\item[\textbf{第三步（求角）：}] 由 $\cos\theta=\frac12$，得 $\theta=60^\circ$（或 $\frac{\pi}{3}$）。
			\par\textbf{理由：}
			特殊角的余弦值。
	\end{description}

	\textbf{\textcolor{CExample}{关键结论：}}
	\textbf{$\cos\theta=\frac12$，$\theta=60^\circ$。}

	\medskip
	\textbf{\textcolor{CExample}{例 3（综合应用）}}

	\textbf{\textcolor{CExample}{题目：}}
	棱长均为 $1$ 的正四面体 $ABCD$ 中，求平面 $ABC$ 与平面 $ABD$ 所成二面角的余弦值。

	\textbf{\textcolor{CExample}{解题步骤：}}
	\begin{description}[style=nextline, leftmargin=2.8em, labelsep=0.5em, itemsep=0.45em]
		\item[\textbf{第一步（确定棱与平面角）：}] 二面角的棱为 $AB$。取 $AB$ 中点 $E$，连接 $CE$、$DE$。由正四面体，$\triangle ABC$ 和 $\triangle ABD$ 均为等边三角形，故 $CE\perp AB$，$DE\perp AB$。因此 $\angle CED$ 是二面角 $C\text{-}AB\text{-}D$ 的平面角，设 $\theta = \angle CED$。
			\par\textbf{理由：}
			定义法要求两条射线均垂直于棱且分别在两个半平面内。

		\item[\textbf{第二步（计算所需边长）：}] 在边长为 $1$ 的等边三角形中，中线长 $\frac{\sqrt3}{2}$，故
			\[
				CE = DE = \frac{\sqrt3}{2},\quad CD = 1.
			\]
			\par\textbf{理由：}
			正四面体所有棱均相等，$CD$ 是侧棱长度 $1$。

		\item[\textbf{第三步（余弦定理求角）：}] 在 $\triangle CED$ 中，由余弦定理得
			\[
				\begin{aligned}
					\cos\theta &= \frac{CE^2+DE^2-CD^2}{2\cdot CE\cdot DE} \\
					&= \frac{\left(\frac{\sqrt3}{2}\right)^2+\left(\frac{\sqrt3}{2}\right)^2-1^2}{2\cdot\frac{\sqrt3}{2}\cdot\frac{\sqrt3}{2}} \\
					&= \frac{\frac34+\frac34-1}{2\cdot\frac34} \\
					&= \frac{\frac12}{\frac32} \\
					&= \frac13.
			\end{aligned}
			\]
			\par\textbf{理由：}
			代入数值并逐步化简；最后一步为分数除法。
	\end{description}

	\textbf{\textcolor{CExample}{关键结论：}}
	\textbf{二面角的余弦值为 $\frac13$。}
\end{examplebox}
